\documentclass[UTF8,fontset=macnew,11pt]{ctexart}

\usepackage[a4paper,margin=2.25cm]{geometry}
\usepackage{amsmath,amssymb}
\usepackage{booktabs}
\usepackage{tabularx}
\usepackage{array}
\usepackage{float}
\usepackage{xcolor}
\usepackage{hyperref}

\hypersetup{
  colorlinks=true,
  linkcolor=blue!50!black,
  citecolor=blue!50!black,
  urlcolor=blue!50!black
}

\setlength{\parskip}{0.34em}
\setlength{\parindent}{2em}
\setlength{\emergencystretch}{2em}

\newcommand{\method}{Resource-NMCTS}
\newcommand{\screen}{Boolean-ring screen}
\newcommand{\anf}{\ensuremath{\mathrm{ANF}}}
\newcommand{\fprm}{\ensuremath{\mathrm{FPRM}}}
\newcommand{\mcts}{\ensuremath{\mathrm{MCTS}}}
\newcommand{\xor}{\ensuremath{\oplus}}
\newcommand{\score}{\ensuremath{\mathrm{score}}}
\newcommand{\ttable}{truth-table}
\newcommand{\mct}{\ensuremath{\mathrm{MCT}}}

\title{面向资源约束量子布尔函数综合的神经蒙特卡洛树搜索方法}
\author{中文论文稿 v19：large frontier policy 版}
\date{2026年7月8日}

\begin{document}
\pagestyle{plain}
\maketitle

\begin{abstract}
量子布尔函数 oracle 综合把经典谓词嵌入量子算法，是 Grover 搜索、量子密码分析和约束求解类量子程序中的基础编译步骤。现有路线通常围绕固定布尔表示、可逆逻辑模板或 CNOT 导向覆盖结构展开，但容错量子计算中的逻辑开销同时受 T-count、CNOT、深度、门数和辅助比特制约。本文将该任务重新表述为代数正规形（algebraic normal form, \anf{}）和固定极性 Reed--Muller（fixed-polarity Reed--Muller, \fprm{}）项集合上的资源约束搜索问题，提出结合神经动作先验、蒙特卡洛树搜索、布尔环线性因子筛选、结构深度策略、depth-frontier policy 和保守 guard 的 \method{} 框架。本文方法不依赖 SSHR 的 parallelotope 候选空间；SSHR 仅作为 small-scale CNOT-oriented baseline 参与比较。实验显示，在 $n\leq6$ 的 177 个传统函数上，Pareto-\method{} 相对 ESOP cube beam 获得 174/0/3 的 score 胜/负/平，平均 score 降低 36.09\%；相对 weighted ESOP-MILP 获得 167/3/7，平均 score 降低 29.84\%。新增 ROS-style LUT proxy 对 $n=3$--$6$ 与 $n=14,15,16,18$ 共 309 个函数运行 ABC $K=3,4,5$ LUT sweep，927/927 行均通过 truth-table 检查；best-$K$ proxy 相对固定 $K=4$ 平均 score 降低 18.12\%，而 \method{} 相对该更强 proxy 仍获得 309/0/0，平均 score 降低 83.77\%。在 $n=18$ 函数级切片上，完整 \method{} 相对 adaptive depth-2 \screen{} 进一步降低 23.85\% 平均 score；在 $n=20$ 边界切片上，adaptive depth-2 \screen{} 相对 AND-direct ANF 降低 11.80\% 平均 score。项集级 $n=20,22,24,28,32,36,40$ scale 实验表明，结构深度策略可以在接近 all-depth adaptive 质量的同时减少约三分之一时间；depth-frontier policy 在 $n=20,28,40$ 上相对 fixed depth-2 获得 35/0/37，平均 score 降低 2.19\%，相对 all-depth depth$\leq4$ 节省 58.69\% 时间。独立 seed 的 $n=24,28,32,40$ 泛化实验中，原 frontier policy 相对 fixed depth-2 为 40/0/56、平均 score 降低 1.85\%；扩大训练集后的 large frontier policy 在同一泛化集上提升到 56/0/40、平均 score 降低 2.34\%，相对旧 policy 为 17/0/79、平均 score 再降 0.49\%，相对 all-depth 仅高 0.10\% 且节省 53.50\% 时间。所有 5640 个项集级方法行均通过 factor-plan \anf{} 展开验证和 emitted X/CNOT/\mct{} circuit 的 GF(2) 多项式符号模拟验证；进一步的 $n=21,22,23$ bridge 与 large-policy $n=23$ rerun 中，240/240 个方法行同时通过完整 \ttable{} oracle 验证、plan 验证和 emitted-circuit 符号验证。新增 emitted-circuit schedule proxy 显示，large frontier policy 在 $n=23$ 完整 bridge 上相对旧 policy 为 1/0/5，平均 score 再降 0.48\%，T-depth proxy 再降 0.45\%。本文的阶段性结论是：\method{} 已形成一条独立于 SSHR、以低 T-count 和低加权逻辑资源为目标的 Boolean oracle synthesis 主线；当前最需要补强的是官方 ROS/RevKit/mockturtle 复现，以及真正的后端映射评估。
\end{abstract}

\noindent\textbf{关键词：}量子 oracle 综合；布尔函数综合；强化学习；蒙特卡洛树搜索；ANF；FPRM；布尔环；资源约束编译

\section*{作者想法整理与当前论文定位}

本稿首先服务于一个清晰的研究主题：基于强化学习与蒙特卡洛树搜索的量子布尔函数综合方法。这个主题的本质不是复现某个已有综合算法，而是把 Boolean oracle synthesis 看作搜索问题：给定布尔函数的 \anf{} 或 \fprm{} 项集合，系统需要在大量可计算、可反算、可验证的代数分解动作中选择一条低资源线路生成路径。强化学习、神经排序和 \mcts{} 的作用，是在候选动作过多时更早找到有资源收益的分解，而不是直接生成不可解释的量子门序列。

因此，本文不应从 SSHR 入手。SSHR 的 parallelotope 结构适合做 small-scale CNOT-oriented synthesis，适合作为重要 baseline，但它不是本文方法的依赖，也不应成为论文主线。更合适的主线是：\anf{}/\fprm{} 项集合搜索定义问题，Boolean-ring screen 扩展可用结构，神经动作先验和 \mcts{} 处理动作组合，depth policy 与 frontier policy 处理搜索预算，guard 保证学习策略不会轻易劣于确定性 baseline。这样的叙事可以自然比较 T-count、CNOT、门数、线路深度和辅助比特，而不是被单个 CNOT 指标锁死。

从投稿角度看，本文当前已有“明显提升”的证据，但证据形态必须如实呈现。小规模传统函数上，本文在低 T-count 和加权 score 上相对 ESOP、weighted ESOP-MILP、direct ANF 和 AND-direct ANF 有稳定优势；高维部分已经从单纯资源估计推进到项集级 plan 验证、emitted-circuit GF(2) 符号验证和 $n=21,22,23$ 完整 \ttable{} bridge。与此同时，本文还不能声称硬件后端最优，也不能声称 CNOT-only 全面超过 SSHR。当前稿件应把优势写成“资源约束逻辑层搜索”，把边界写成“仍需外部工具链和后端评估”。

下一步最值得投入的方向也很明确。第一，在当前 ROS-style LUT proxy 的基础上继续接入或复现官方 ROS、mockturtle/RevKit 风格 reversible-toolchain，把比较从 proxy 扩展到更标准的外部流程。第二，large frontier policy 已把独立泛化集上相对 all-depth 的平均 score gap 从 0.61\% 压到 0.10\%，但代价是相对旧 policy 的时间增加，因此仍需学习更好的质量-时间-辅助线生命周期联合目标。第三，在当前 schedule proxy 的基础上进一步做真实后端映射或至少更精细的逻辑层调度，例如 T-depth 模板、辅助量子位生命周期、反算区间和可能的 routing 压力。完成这些以后，论文就更接近“方法贡献”而不是“阶段性项目报告”。

\section{引言}

给定布尔函数 $f:\{0,1\}^n\rightarrow\{0,1\}$，量子 oracle 通常实现为
\begin{equation}
  |x\rangle |y\rangle \mapsto |x\rangle |y\xor f(x)\rangle .
  \label{eq:oracle}
\end{equation}
这一步看似只是把经典逻辑接入量子程序，但在容错量子计算中，它往往决定了非 Clifford 资源和辅助比特压力。若直接按 \anf{} 单项式逐项实现，线路语义透明，但容易重复计算相似结构；若只优化 CNOT 或只压缩某一种图表示，又可能把成本转移到 T-count、深度或辅助比特上。因此，本文关注的问题不是“如何复现某个已有 CNOT 优化算法”，而是“如何把 Boolean oracle synthesis 写成一个可由 AI 辅助的多资源搜索问题”。

本文采用 \anf{} 作为语义基准：
\begin{equation}
  f(x)=\bigoplus_{m\in\mathcal{M}}\prod_{i\in m}x_i ,
  \label{eq:anf}
\end{equation}
其中 $\mathcal{M}$ 是需要综合的 square-free monomial 集合。该表示有两个优势。第一，任一候选线路都可以回到 GF(2) 多项式进行等价验证。第二，搜索动作可以直接定义为对项集合的代数变换，例如公共因子提取、\fprm{} 极性重写、线性奇偶因子筛选和递归子问题分解。它的困难也很明确：原始 \anf{} 不显式给出共享结构，动作空间随项数和变量数迅速增长，必须用搜索策略、神经先验和保守兜底共同控制。

本文的一句话论证是：在逻辑层量子布尔函数 oracle 综合中，把 \anf{}/\fprm{} 项集合当作搜索状态，并用资源感知神经搜索选择可计算、可反算的代数分解动作，可以在传统函数和高维项集合上降低 T-count 与加权逻辑资源；其边界是当前结果仍停留在逻辑层，不包含 routing、magic-state factory、真实噪声模型或硬件后端映射。

本文贡献概括如下：
\begin{itemize}
  \item 提出独立于 SSHR 的 \method{} 搜索表述，以 \anf{}/\fprm{} 项集合为状态，以逻辑资源 score 为目标，把 Boolean oracle synthesis 写成可验证的组合搜索问题。
  \item 引入 \screen{}，在 square-free Boolean ring 中允许线性因子与 quotient 共享变量，从而扩大高维可搜索因子结构。
  \item 将 AI 模块分为动作排序、结构深度选择、depth-frontier policy 和 baseline-preserving guard，避免把学习策略写成不可控的黑箱生成器。
  \item 在 $n\leq6$ 传统函数、$n=18,20$ 函数级边界、$n=20$ 到 $40$ 项集级 scale、$n=21,22,23$ 完整 \ttable{} bridge 以及 emitted-circuit schedule proxy 上给出资源、正确性和逻辑层时序证据。
\end{itemize}

\section{相关工作与本文定位}

\subsection{布尔表示驱动的 oracle 综合}

Xor-And-Inverter Graphs（XAG）把 XOR 和 AND 结构分开，使 AND 节点数量与低 T-count oracle 编译直接相关。Meuli 等人讨论了 multiplicative complexity 在低 T-count oracle 编译中的作用\cite{meuli2019multiplicative}，并将 XAG 用于量子编译\cite{meuli2022xag}。ROS 将 oracle synthesis 表述为 resource-constrained LUT mapping\cite{meuli2020ros}，进一步说明布尔表示选择会影响 T-count、辅助比特和线路深度。BDD-based reversible synthesis 则展示了面向大函数的可扩展布尔表示路径\cite{wille2009bdd}。后端感知 oracle synthesis 开始把逻辑综合和 fault-tolerant back-end 指标连接起来\cite{yu2025backend}。

本文与上述工作的共同点是承认布尔中间表示决定逻辑资源；差异在于本文不固定使用 XAG、LUT 或 BDD，而是在 \anf{}/\fprm{} 项集合上显式搜索可反算的因子结构。这样做的优点是语义验证直接，代价是候选动作空间更大，需要 AI 和 guard 控制搜索。

SSHR 利用 parallelotope 的空间结构做 small-scale Boolean function 的 CNOT-oriented synthesis\cite{zheng2025sshr}。本文不使用 SSHR 的候选空间，也不把论文目标设为 CNOT-only 最优。SSHR 在本文中是一个重要但定位明确的 baseline：它帮助说明本文在低 T-count 和加权 score 上的优势，也提醒我们在 CNOT 指标上必须诚实报告 trade-off。

\subsection{学习式搜索与线路优化}

学习式搜索已经出现在经典布尔电路设计、量子线路综合和 ZX-calculus 优化中。ShortCircuit 使用 AlphaZero 风格搜索进行经典布尔电路设计\cite{tsaras2024shortcircuit}；Gumbel AlphaZero 被用于 Clifford+T unitary synthesis\cite{rietsch2024unitary}；强化学习和图神经网络也被用于 ZX-calculus rewrite-rule 选择\cite{riu2025rlzx}；MonteQ 使用 \mcts{} 处理量子线路综合中的动作排序问题\cite{machiya2026monteq}。这些工作证明，离散线路综合中的动作选择可以受益于学习策略。

本文的区别在于状态不是任意量子门序列，也不是 ZX 图，而是 Boolean oracle 的 \anf{}/\fprm{} 项集合。神经模型不直接输出线路，而是排序或选择候选代数分解动作；最终线路仍由可验证的 compute/action/uncompute 结构生成。这一设计降低了 AI 生成错误线路的风险，也使每个候选都可以经过 \anf{} 展开、GF(2) 符号模拟或完整 \ttable{} 检查。

\section{问题定义与资源模型}

本文输入是布尔函数的 \ttable{} 或 \anf{} 项集合，输出是实现式~\eqref{eq:oracle} 的逻辑层可逆线路。线路由 X、CNOT 和多控 Toffoli（\mct{}）抽象门组成，并在资源统计中估计 T-count、CNOT、depth、gate count 和 peak ancilla。搜索排序使用统一分数
\begin{equation}
  \score = T + 0.04\,\mathrm{CNOT} + 0.015\,\mathrm{depth}
         + 0.01\,\mathrm{gates} + 2\,\mathrm{ancilla}.
  \label{eq:score}
\end{equation}
该分数不是物理设备的时间或错误率模型，而是逻辑层候选排序目标。为避免单一分数掩盖 trade-off，实验同时报告各个资源分量。

正确性验证分三层。第一，在函数级小规模和边界切片中构造完整 \ttable{}，逐点检查 oracle 输出。第二，在大规模项集实验中，将 factor plan 展开回 \anf{} 项集合，检查目标多项式是否等价。第三，对 emitted X/CNOT/\mct{} circuit 做 GF(2) 多项式符号模拟，检查输出线多项式、输入线复原和辅助线清零。第三层验证不等同于 $2^n$ 个输入的完整枚举，但它检查的是最终线路而不只是搜索 plan，因而强于只报告资源估计。

\begin{table}[H]
\centering
\small
\caption{本文核心术语。}
\label{tab:terms}
\begin{tabularx}{\linewidth}{>{\raggedright\arraybackslash}p{0.27\linewidth}>{\raggedright\arraybackslash}X}
\toprule
术语 & 本文含义 \\
\midrule
\method{} & 在 \anf{}/\fprm{} 项集合上进行资源加权搜索的主框架，包括神经动作先验、\mcts{}、guard 和 portfolio 选择。 \\
\screen{} & 在 square-free Boolean ring 中搜索可复用线性因子的结构筛选分支，允许线性因子与 quotient 共享变量。 \\
Depth policy & 在 single、depth-1 和 depth-2 screen 之间选择的结构级神经策略。 \\
Depth-frontier policy & 在 depth-2、depth-3 和 depth-4 screen 之间选择的高预算质量-时间折中策略。 \\
Guard & 以确定性 baseline 为兜底的保守门控，只有当学习候选不劣于 baseline 时才接受。 \\
Truth-table bridge & 在 $n=21,22,23$ 构造完整 \ttable{}，把大规模项集符号验证和函数级完整验证连接起来。 \\
SSHR & CNOT-oriented small-scale baseline；本文方法不使用 SSHR parallelotope 候选。 \\
\bottomrule
\end{tabularx}
\end{table}

\section{方法}

\subsection{搜索状态、动作与线路生成}

\method{} 的状态是当前尚未综合的项集合 $\mathcal{M}$。一次动作选择一个可计算、可反算的代数结构，将 $\mathcal{M}$ 分解为 quotient 子问题和 rest 子问题，再递归生成 compute/action/uncompute 形式的线路。候选动作包括公共单项式因子、\fprm{} 极性重写后的公共项、线性奇偶因子、布尔环线性因子，以及由神经模型扩展的 root actions。

动作 $a$ 被接受后，框架会估计其即时资源收益、子问题大小、辅助比特峰值和后续可分解性。神经动作先验对候选排序，\mcts{} 在有限预算内探索候选组合，rollout 使用确定性资源估计和已有 baseline 候选。最终 portfolio 在多个候选线路中按式~\eqref{eq:score} 选取最优者。该过程的关键不是让 AI 直接创造量子门，而是让 AI 在大量可验证的代数动作中更早访问有资源收益的分支。

\subsection{布尔环线性因子}

传统线性因子筛选容易要求线性因子变量和 quotient 变量不相交。这个限制便于实现，但会排除 square-free Boolean ring 中合法的重叠结构。本文的 \screen{} 允许二者共享变量，并在展开时使用 $x_i^2=x_i$ 与 GF(2) 偶数项消去。例如
\begin{equation}
  (x_i\xor x_j)x_i m = x_i m \xor x_i x_j m .
  \label{eq:boolean-ring}
\end{equation}
线路层面仍可用 CNOT 计算线性奇偶量，再以该辅助量控制 quotient 子线路。这样，\screen{} 不是简单增加 beam 宽度，而是把原先被实现约束排除的代数候选重新放回可综合空间。

\subsection{神经策略、结构策略与保守 guard}

本文把 AI 的角色拆成三类。第一类是动作排序：模型根据候选覆盖率、项数变化、变量覆盖密度、次数分布和即时资源估计对候选动作打分。第二类是结构选择：depth policy 判断当前项集合适合 single、depth-1 还是 depth-2 \screen{}。第三类是 frontier 选择：depth-frontier policy 在 depth-2、depth-3 和 depth-4 \screen{} 之间选择，目标是在接近 depth-4 或 all-depth frontier 的质量时减少全深度枚举开销。

保守 guard 用来限制学习策略的风险。若学习候选相对确定性 baseline 出现 score 退化，则返回 baseline。Depth-2 skip guard 采用 zero-false-skip 阈值，优先避免把应该运行 depth-2 的项集合错误跳过。Screen-gate 则用于判断在某些边界函数上是否可以跳过昂贵的 Resource tail。这样的设计降低了策略模型分布外失误对资源结果的影响，但也意味着部分 AI 贡献表现为运行时间节省或质量-时间折中，而不是无条件的质量支配。

\begin{table}[H]
\centering
\small
\caption{\method{} 的模块与证据钩子。}
\label{tab:modules}
\begin{tabularx}{\linewidth}{>{\raggedright\arraybackslash}p{0.24\linewidth}>{\raggedright\arraybackslash}X>{\raggedright\arraybackslash}p{0.25\linewidth}}
\toprule
模块 & 机制 & 主要证据类型 \\
\midrule
ANF/FPRM 搜索状态 & 把 oracle 综合写成项集合分解问题，所有候选可回到 GF(2) 多项式验证。 & 小规模 truth-table 验证与项集级 ANF 验证。 \\
Boolean-ring screen & 搜索允许共享变量的线性因子，递归处理 quotient/rest 子问题。 & $n=18,20$ 函数级边界和 $n=20$--$40$ scale。 \\
神经动作先验与 \mcts{} & 对候选动作排序并在有限预算内组合分解动作。 & 传统函数上的搜索消融和 portfolio 改善。 \\
Depth policy / guard & 学习 screen 深度或跳过较深搜索，并用 baseline-preserving 规则兜底。 & held-out 项集时间节省和无 score-loss guard。 \\
Depth-frontier policy & 学习高预算 depth-2/3/4 质量前沿的折中点。 & $n=20,28,40$ scale 与 $n=21,22,23$ bridge。 \\
\bottomrule
\end{tabularx}
\end{table}

\section{实验设计}

实验分为六个层次。第一层是 $n\leq6$ 传统函数，用 direct ANF、AND-direct ANF、ESOP cube beam、weighted ESOP-MILP、ABC/BDD 估计、ROS-style LUT proxy 和 SSHR-H/SSHR-I 作为 baseline，主要回答小规模函数上是否形成低 T-count 和低加权资源优势。第二层是 $n=14,15,16,18$ 的高维外部导出集，对 ABC-AIG/XAG/LUT/BDD 和 ROS-style LUT proxy 做统一 truth-table 校验，主要回答外部布尔表示路线在高维上是否会削弱本文优势。第三层是 $n=18,20$ 函数级边界，完整构造 \ttable{} 并验证 emitted oracle circuit，主要回答高维函数上哪些搜索分支仍可运行。第四层是 $n=20$ 到 $40$ 项集级 scale，不构造完整 \ttable{}，而是对 plan 和 emitted circuit 做符号等价验证，主要回答结构策略是否可扩展，并通过独立 seed 的 $n=24,28,32,40$ 泛化集检查 frontier policy 的稳定性。第五层是 $n=21,22,23$ truth-table bridge，对一小组生成式 \anf{} 函数重新构造完整 \ttable{}，专门检查项集级符号验证是否与函数级 oracle 验证一致。第六层是 emitted-circuit schedule proxy，在不做硬件 mapping 的前提下估计并行逻辑深度、CNOT-depth proxy、T-depth proxy 和显式辅助线 lifetime area。

所有 paired comparison 只纳入函数名或项集名匹配、语义正确且未 timeout 的结果。Timeout 行保留为运行边界证据，但不参与正确结果均值比较。所有结论限定在逻辑层资源，不声称硬件 mapping、连通性约束或噪声模型下最优。

\section{结果}

\subsection{传统函数上的低 T-count 与低 score 优势}

表~\ref{tab:traditional} 给出 $n\leq6$ 传统函数上的平均资源。相对 direct ANF，Pareto-\method{} 平均 T-count 降低 72.25\%，平均 score 降低 67.80\%。相对 ESOP cube beam，Pareto-\method{} 获得 174/0/3 的 score 胜/负/平，平均 score 降低 36.09\%。相对 weighted ESOP-MILP，获得 167/3/7，平均 score 降低 29.84\%。SSHR-H 的平均 CNOT 最低，因此本文不能主张 CNOT-only 支配；更准确的结论是，本文用一定 CNOT 与辅助比特 trade-off 换取更低 T-count 和更低加权 score。

\begin{table}[H]
\centering
\small
\caption{$n\leq6$ 传统函数切片上的平均逻辑资源。}
\label{tab:traditional}
\begin{tabular}{lrrrrr}
\toprule
方法 & 平均 T & 平均 CNOT & 平均 depth & 平均 ancilla & 平均 score \\
\midrule
Direct ANF & 184.1 & 134.2 & 134.6 & 1.33 & 194.3 \\
AND-direct ANF & 101.0 & 166.5 & 166.9 & 2.18 & 115.2 \\
\method{} & 43.9 & 89.9 & 94.5 & 1.92 & 53.2 \\
Pareto-\method{} & \textbf{40.4} & 83.0 & \textbf{88.0} & 2.03 & \textbf{49.6} \\
ESOP-MILP & 83.6 & 133.5 & 159.6 & 2.32 & 96.7 \\
SSHR-H & 81.0 & \textbf{67.6} & 88.7 & 1.36 & 88.2 \\
\bottomrule
\end{tabular}
\end{table}

\subsection{ROS-style LUT proxy 外部对比}

官方 ROS 把 oracle synthesis 写成 resource-constrained LUT mapping，并配合垃圾管理策略降低辅助线压力。当前环境尚不能直接运行官方 ROS/RevKit/mockturtle 流程，因此本文新增一个明确标注的 ROS-style LUT proxy：对同一导出 BLIF benchmark 运行 ABC `if -K`，在 $K=3,4,5$ 之间做 sweep；每个映射后的 LUT 网络按 local \anf{} compute/action/uncompute 方式估计 X/CNOT/\mct{} 逻辑资源，再按式~\eqref{eq:score} 选择 best-$K$。这不是官方 ROS 复现，也没有 SAT garbage management，但它比单个固定 $K=4$ 的 ABC-LUT baseline 更强，更适合作为 LUT 路线的压力测试。

该 proxy 覆盖 $n=3$--$6$ 的 177 个传统函数，以及 $n=14,15,16,18$ 的 132 个高维导出函数，共 309 个函数。三档 $K$ sweep 产生 927 个外部映射行，全部通过 truth-table 检查，没有 timeout、error 或 incorrect 行。Best-$K$ proxy 相对固定 $K=4$ 自身获得 219/0/90 的 score 胜/负/平，平均 score 降低 18.12\%；因此它确实是比原 ABC-LUT 更强的 LUT baseline。在这个更强 baseline 上，\method{} 仍在 309/309 个函数上取得更低 score，平均 score 降低 83.77\%；Pareto-\method{} 为 309/0/0、平均 score 降低 84.27\%。高维 profile variant 在 132 个高维函数上相对该 proxy 为 132/0/0、平均 score 降低 97.38\%。

\begin{table}[H]
\centering
\scriptsize
\caption{LUT proxy 外部对比。该 proxy 不是官方 ROS/RevKit/mockturtle 复现。}
\label{tab:ros-lut}
\begin{tabular}{llrrr}
\toprule
比较 & 指标 & 函数数 & 胜/负/平 & 平均 $\Delta$ \\
\midrule
\method{} vs ROS-style LUT proxy & score & 309 & 309/0/0 & $-83.77\%$ \\
Profile-\method{} vs ROS-style LUT proxy & score & 132 & 132/0/0 & $-97.38\%$ \\
Pareto-\method{} vs ROS-style LUT proxy & score & 309 & 309/0/0 & $-84.27\%$ \\
AND-direct ANF vs ROS-style LUT proxy & score & 309 & 307/2/0 & $-67.45\%$ \\
ROS-style LUT proxy vs fixed $K=4$ LUT & score & 309 & 219/0/90 & $-18.12\%$ \\
\bottomrule
\end{tabular}
\end{table}

\subsection{高维函数级边界}

在 $n=18$ 的 12 个随机 \anf{} 函数上，adaptive depth-2 \screen{} 相对 single screen 获得 11/0/1，平均 score 从 11349.06 降至 10670.94；完整 \method{} 平均 score 为 7315.62，相对 adaptive screen 进一步降低 23.85\%。这说明中等高维上 screen 是强候选生成器，但不能完全替代 Resource tail。

在 $n=20$ 边界函数上，FPRM root beam 和 fast linear-pair 均触发 300 s task timeout。Adaptive depth-2 screen 相对 single screen 获得 5/0/1，平均 score 降低 7.47\%；集成该分支后，\method{} 相对 AND-direct ANF 获得 5/0/1，平均 score 降低 11.80\%。该切片中完整 Resource tail 与 adaptive screen 资源持平，因此 screen-gated \method{} 可把平均时间从 77.10 s 降至 20.71 s。这个结果支持运行时门控，但也提示 $n=20$ 切片的主要质量收益来自结构筛选而非深层 \mcts{} tail。

\begin{table}[H]
\centering
\scriptsize
\caption{$n=18,20$ 函数级边界结果。}
\label{tab:highdim}
\begin{tabular}{llrrrr}
\toprule
切片 & 方法 & 平均 T & 平均 CNOT & 平均 score & 平均时间(s) \\
\midrule
$n=18$ & Single screen & 10389.33 & 16268.42 & 11349.06 & 0.82 \\
$n=18$ & Adaptive depth-2 screen & 9767.00 & 15301.58 & 10670.94 & 4.39 \\
$n=18$ & \method{} & \textbf{6639.33} & \textbf{11380.92} & \textbf{7315.62} & 226.18 \\
$n=20$ & AND-direct ANF & 21516.67 & 33635.67 & 23491.15 & 10.41 \\
$n=20$ & Single screen & 20786.00 & 32492.50 & 22695.15 & 10.70 \\
$n=20$ & Adaptive depth-2 screen & \textbf{19535.33} & \textbf{30542.50} & \textbf{21331.24} & 20.67 \\
$n=20$ & \method{} & \textbf{19535.33} & \textbf{30542.50} & \textbf{21331.24} & 77.10 \\
$n=20$ & Screen-gated \method{} & \textbf{19535.33} & \textbf{30542.50} & \textbf{21331.24} & 20.71 \\
\bottomrule
\end{tabular}
\end{table}

\subsection{结构级 AI 与高预算 frontier}

Depth policy 在 $n=14,16,18$ 项集上训练，并在 held-out $n=20$ 项集上评估。Policy 相对 single screen 获得 42/0/6，平均 score 降低 6.57\%；相对 all-depth adaptive 平均 score 只高 0.28\%，但平均运行时间降低 34.71\%。保守 depth-2 skip guard 在 held-out $n=20$ test 中没有 false skip，96 个样本全部与 fixed depth-2 screen score 持平，并相对 all-depth adaptive 节省 24.74\% 平均时间。

为判断 fixed depth-2 是否已经达到质量上限，本文在 $n=20,28,40$ 上运行 depth-3/4 \screen{}。Depth-4 相对 fixed depth-2 获得 49/0/23，平均 score 降低 3.10\%，但运行时间增加 681.55\%。因此，depth-3/4 是高预算质量模式，而不是默认快速模式。Depth-frontier policy 将这一质量前沿学习化：在 $n=20,28,40$ scale harness 中，它相对 fixed depth-2 获得 35/0/37，平均 score 降低 2.19\%；相对 all-depth adaptive depth$\leq4$ 平均 score 只高 0.97\%，但节省 58.69\% 时间。独立 seed 的 $n=24,28,32,40$ 泛化集进一步显示，frontier policy 相对 fixed depth-2 获得 40/0/56，平均 score 降低 1.85\%；相对 fixed depth-4 只高 0.61\% score，并节省 23.40\% plan time。

为了补强“AI 结构策略质量收益仍偏弱”的问题，本文进一步扩大 depth-frontier policy 的训练集：训练样本从 96 增至 192，验证样本从 36 增至 72，held-out 测试样本从 32 增至 48，并把隐藏层宽度从 96 提高到 160。Large frontier policy 在 held-out 测试上相对 oracle depth-2/3/4 frontier 的平均 score gap 从旧模型的 0.80\% 降至 0.04\%，同时相对 all-depth frontier evaluation 节省 51.30\% 时间。在同一独立 seed 的 $n=24,28,32,40$ 泛化集上，large frontier policy 相对 fixed depth-2 获得 56/0/40，平均 score 降低 2.34\%；相对旧 frontier policy 为 17/0/79，平均 score 再降 0.49\%；相对 all-depth adaptive 只高 0.10\%，仍节省 53.50\% 时间。该结果说明结构级 AI 的质量 gap 可以通过更充分的 frontier teacher 数据显著压缩，但它也更频繁选择 depth-3/4，因此比旧 policy 更慢。

\begin{table}[H]
\centering
\scriptsize
\caption{结构级策略与 depth-frontier 结果。}
\label{tab:policy}
\begin{tabular}{llrrr}
\toprule
设置 & 比较 & score 胜/负/平 & 平均 $\Delta$ score & 平均 $\Delta$ time \\
\midrule
held-out $n=20$ & depth policy vs single screen & 42/0/6 & $-6.57\%$ & $+2328.66\%$ \\
held-out $n=20$ & depth policy vs all-depth adaptive & 0/6/42 & $+0.28\%$ & $-34.71\%$ \\
held-out $n=20$ & depth-2 guard vs fixed depth-2 & 0/0/96 & $+0.00\%$ & $-0.54\%$ \\
scale $n=20,28,40$ & screen depth-4 vs fixed depth-2 & 49/0/23 & $-3.10\%$ & $+681.55\%$ \\
scale $n=20,28,40$ & frontier policy vs fixed depth-2 & 35/0/37 & $-2.19\%$ & $+503.20\%$ \\
scale $n=20,28,40$ & frontier policy vs all-depth & 0/16/56 & $+0.97\%$ & $-58.69\%$ \\
generalization $n=24,28,32,40$ & frontier policy vs fixed depth-2 & 40/0/56 & $-1.85\%$ & $+456.43\%$ \\
generalization $n=24,28,32,40$ & frontier policy vs depth-4 & 0/19/77 & $+0.61\%$ & $-23.40\%$ \\
large held-out & large frontier vs oracle frontier & 0/3/45 & $+0.04\%$ & $-51.30\%$ \\
large generalization & large frontier vs fixed depth-2 & 56/0/40 & $-2.34\%$ & $+563.80\%$ \\
large generalization & large frontier vs old frontier & 17/0/79 & $-0.49\%$ & $+119.26\%$ \\
large generalization & large frontier vs all-depth & 0/6/90 & $+0.10\%$ & $-53.50\%$ \\
\bottomrule
\end{tabular}
\end{table}

\subsection{项集级 scale 与完整验证桥接}

项集级 scale 覆盖 $n=20,22,24,28,32,36,40$，主 scale、extended scale、depth-frontier、frontier-policy rerun、独立泛化 run 和 large frontier policy 泛化 run 合计 5640 个方法行。所有方法行均通过两层验证：factor plan 的 \anf{} 展开验证为 5640/5640，emitted X/CNOT/\mct{} circuit 的 GF(2) 多项式符号模拟验证为 5640/5640，mismatch 总数为 0。该结果说明大规模项集资源表不是只统计 plan，但仍需承认它不是完整 \ttable{} 枚举。

Truth-table bridge 进一步把完整验证推进到 $n=21,22,23$。在 18 个生成式 \anf{} 函数和 10 个方法上，180/180 个方法行通过完整 \ttable{} oracle 验证、plan \anf{} 展开验证和 emitted-circuit GF(2) 符号验证，mismatch 为 0。Large frontier policy 又在相同 seed 的 $n=23$ 六函数切片上重新运行 60 个方法行，也全部通过完整 \ttable{} oracle、plan 和 emitted-circuit 验证。因此，当前 bridge 证据合计为 240/240 个方法行通过三层验证。其中原 $n=23$ bridge 平均 truth-table 构造时间为 209.54 s/function，large rerun 为 208.61 s/function。该 bridge 还复现了 frontier 信号：原 frontier policy 在 $n=23$ 上相对 fixed depth-2 获得 4/0/2，平均 score 降低 1.88\%；large frontier policy 提升为 5/0/1，平均 score 降低 2.36\%。相对旧 frontier policy，large frontier 在 $n=23$ 上为 1/0/5，平均 score 再降 0.48\%，T-depth proxy 再降 0.45\%。

\begin{table}[H]
\centering
\scriptsize
\caption{大规模项集验证与 truth-table bridge。}
\label{tab:verification}
\begin{tabularx}{\linewidth}{>{\raggedright\arraybackslash}p{0.35\linewidth}>{\raggedright\arraybackslash}p{0.22\linewidth}>{\raggedright\arraybackslash}X}
\toprule
验证项或比较 & 结果 & 说明 \\
\midrule
项集级 plan \anf{} 验证 & 5640/5640 & 覆盖 $n=20$ 到 $40$ 的 scale/frontier/frontier-policy rerun、独立 seed 泛化 run 和 large frontier policy 泛化 run。 \\
项集级 emitted-circuit 符号验证 & 5640/5640 & GF(2) 多项式模拟通过，mismatch 总数为 0。 \\
$n=21,22,23$ 与 large $n=23$ 完整 \ttable{} oracle 验证 & 240/240 & 18 个原 bridge 函数加 6 个 large-policy n=23 rerun 函数，所有 emitted circuit 逐点验证通过。 \\
$n=21,22$ frontier policy vs fixed depth-2 & 8/0/4 & 平均 score $-3.50\%$，plan time $+634.71\%$。 \\
$n=23$ frontier policy vs fixed depth-2 & 4/0/2 & 平均 score $-1.88\%$，plan time $+759.95\%$。 \\
$n=23$ frontier policy vs all-depth & 0/1/5 & 平均 score $+0.61\%$，plan time $-48.77\%$。 \\
large $n=23$ frontier vs fixed depth-2 & 5/0/1 & 平均 score $-2.36\%$，plan time $+790.62\%$。 \\
large $n=23$ frontier vs old frontier & 1/0/5 & 平均 score $-0.48\%$，T-depth proxy $-0.45\%$。 \\
\bottomrule
\end{tabularx}
\end{table}

\subsection{逻辑层 schedule proxy 与辅助线生命周期}

为避免只报告静态资源总量，本文进一步在 emitted X/CNOT/\mct{} oracle circuit 上计算逻辑层 schedule proxy。该分析不做硬件 mapping，也不使用连通性、routing、噪声或 magic-state factory 模型；它只把每个 emitted gate 按线冲突做并行层调度，并用 MCT 分解的 CNOT 数和 T-stage proxy 估计 CNOT-depth 与 T-depth，同时记录显式辅助线从第一次使用到最后一次使用的 lifetime area。因此，这组指标应理解为后端相关的逻辑层代理指标，而不是物理设备运行时间。

表~\ref{tab:schedule} 显示，depth-frontier policy 的质量收益在 T-depth proxy 上也成立。在独立 seed 的 $n=24,28,32,40$ 泛化集中，frontier policy 相对 fixed depth-2 获得 40/0/56 的 score 胜/负/平，平均 score 降低 1.85\%；同一比较下 T-depth proxy 也为 40/0/56，平均降低 1.85\%。Large frontier policy 在同一泛化集上把 score 提升为 56/0/40、平均降低 2.34\%，T-depth proxy 平均降低 2.28\%，但显式辅助线 lifetime area 增幅扩大到 26.13\%。在 $n=21,22$ 完整 \ttable{} bridge 中，frontier policy 相对 fixed depth-2 的 T-depth proxy 为 8/0/4、平均降低 3.32\%；在新增 $n=23$ bridge 中为 4/0/2、平均降低 1.69\%。Large frontier policy 在 $n=23$ bridge 中相对 fixed depth-2 的 T-depth proxy 平均降低 2.14\%，相对旧 policy 再降低 0.45\%，但辅助线 lifetime area 相对旧 policy 增加 4.00\%。因此，新策略应写成质量增强型 frontier policy，而不是纯省时策略。

\begin{table}[H]
\centering
\tiny
\caption{逻辑层 schedule proxy 结果。所有指标来自 emitted circuit，不包含硬件 mapping。}
\label{tab:schedule}
\begin{tabular}{@{}p{0.20\linewidth}p{0.20\linewidth}p{0.14\linewidth}p{0.15\linewidth}p{0.10\linewidth}@{}}
\toprule
设置 & 对比 & score & T-depth proxy & aux area \\
\midrule
$n=24,28,32,40$ schedule 泛化 & frontier policy vs fixed depth-2 & 40/0/56, $-1.85\%$ & 40/0/56, $-1.85\%$ & $+20.09\%$ \\
$n=24,28,32,40$ schedule 泛化 & frontier policy vs depth-4 & 0/19/77, $+0.61\%$ & 0/19/77, $+0.55\%$ & $-5.42\%$ \\
$n=21,22$ schedule bridge & frontier policy vs fixed depth-2 & 8/0/4, $-3.50\%$ & 8/0/4, $-3.32\%$ & $+32.93\%$ \\
$n=21,22$ schedule bridge & frontier policy vs depth-4 & 0/2/10, $+0.32\%$ & 0/2/10, $+0.32\%$ & $-4.01\%$ \\
$n=23$ schedule bridge & frontier policy vs fixed depth-2 & 4/0/2, $-1.88\%$ & 4/0/2, $-1.69\%$ & $+29.49\%$ \\
$n=23$ schedule bridge & frontier policy vs depth-4 & 0/1/5, $+0.61\%$ & 0/1/5, $+0.52\%$ & $-5.09\%$ \\
large $n=24,28,32,40$ & large frontier vs fixed depth-2 & 56/0/40, $-2.34\%$ & $-2.28\%$ & $+26.13\%$ \\
large $n=24,28,32,40$ & large frontier vs all-depth & 0/6/90, $+0.10\%$ & $+0.10\%$ & $-1.17\%$ \\
large $n=23$ bridge & large frontier vs old frontier & 1/0/5, $-0.48\%$ & $-0.45\%$ & $+4.00\%$ \\
\bottomrule
\end{tabular}
\end{table}

\section{证据地图与边界}

表~\ref{tab:evidence-map} 将当前可写入论文的主张与证据对应起来。最重要的写法是把“AI”讲成可拆分、可验证的搜索控制模块，而不是泛泛声称使用强化学习；把“SSHR”讲成 CNOT-oriented baseline，而不是本文方法来源；把“高维结果”讲成项集级符号验证加 bridge，而不是所有 $n=40$ 函数都做了完整 \ttable{}。

\begin{table}[H]
\centering
\small
\caption{当前主张、证据与投稿边界。}
\label{tab:evidence-map}
\begin{tabularx}{\linewidth}{>{\raggedright\arraybackslash}p{0.28\linewidth}>{\raggedright\arraybackslash}X>{\raggedright\arraybackslash}p{0.23\linewidth}}
\toprule
主张 & 当前证据 & 边界 \\
\midrule
本文路线独立于 SSHR & 状态、动作和验证都定义在 \anf{}/\fprm{} 项集合上；SSHR 只出现在 baseline 比较中。 & 不能声称 CNOT-only 支配 SSHR。 \\
低 T-count 与低加权 score 是主优势 & $n\leq6$ 传统函数相对 ESOP、weighted ESOP-MILP、direct ANF 和 AND-direct ANF 有稳定优势。 & 资源权重仍是逻辑层人为设定。 \\
\screen{} 是高维有效结构筛选 & $n=18,20$ 函数级切片和 $n=20$--$40$ 项集级 scale 均显示相对 single screen 的 score 改善。 & 高维不是全局最优性证明。 \\
结构级 AI 能降低搜索成本并压缩质量 gap & Depth policy、skip guard、screen-gate、frontier policy 和 large frontier policy 均有 held-out、scale 或 bridge 证据；large policy 在泛化集上相对旧 policy 为 17/0/79。 & Large policy 是质量增强型，时间和辅助线生命周期代价更高。 \\
大规模结果具有语义可信度 & 5640/5640 项集级方法行通过 plan 和 emitted-circuit 符号验证；240/240 bridge 行通过完整 \ttable{} 验证。 & $n=24$--$40$ 仍未做完整 \ttable{} 枚举。 \\
逻辑层后端相关 proxy 可解释 trade-off & schedule 泛化和 bridge 显示 frontier policy 同步降低 T-depth proxy；large policy 进一步降低 T-depth，但增加辅助线 lifetime area。 & 不是硬件 mapping、routing 或噪声模型结论。 \\
投稿前仍需补强外部比较 & 当前已有 ESOP、ABC、BDD、SSHR、ROS-style LUT proxy 对比和 toolchain readiness 审计。 & 官方 ROS、mockturtle、RevKit 流程仍待复现。 \\
\bottomrule
\end{tabularx}
\end{table}

\section{讨论}

本文最稳妥的贡献表述是“资源约束搜索”，不是“AI 全面战胜传统综合”。当前证据显示，\anf{}/\fprm{} 项集合提供了清晰的语义验证基础，\screen{} 提供了高维结构候选，\mcts{} 与神经先验提供了候选组合能力，depth policy 与 guard 控制搜索成本，depth-frontier policy 则把高预算质量前沿部分学习化。Large frontier policy 进一步说明，只要提供更充分的 frontier teacher 数据，结构级 AI 的质量 gap 可以从 0.61\% 量级压到 0.10\% 量级。这些模块合在一起，形成了一条比单纯复现 SSHR 更适合投稿的逻辑层 Boolean oracle synthesis 主线。

同时，本文必须明确三个限制。第一，SSHR 的 CNOT-oriented 优势仍存在；在 CNOT-only 或 CNOT-depth profile 下，本文优势会收窄。第二，高维神经策略的独立质量收益虽已通过 large frontier policy 变得更清楚，但目前主要发生在结构深度选择层，而不是任意动作级神经搜索全面优于手工结构。第三，schedule proxy 只是 emitted-circuit 层的逻辑调度代理，本文仍没有处理硬件 mapping、连通性、routing、magic-state factory 或噪声下可靠性，因而不能直接推出物理设备上更快或更可靠。

下一步最有价值的提升目标不是继续堆内部表格，而是补两个外部缺口。其一，在 ROS-style LUT proxy 之外复现或接入官方 ROS、mockturtle/RevKit 风格 reversible-toolchain，使比较从 proxy 扩展到更标准的综合流程。其二，把 large frontier policy 从质量增强型策略改造成质量-时间-辅助线生命周期联合优化策略，减少相对旧 policy 的时间增幅和相对 fixed depth-2 的 lifetime area 增幅。若这两点完成，论文主张会从“已有一条合理路线”提升到“有更强外部说服力的算法贡献”。

\section{结论}

本文提出 \method{}，将量子布尔函数 oracle 综合建模为 \anf{}/\fprm{} 项集合上的资源约束搜索问题，并结合神经动作先验、\mcts{}、布尔环线性因子筛选、结构深度策略、depth-frontier policy 和保守 guard 生成逻辑层可逆线路。实验表明，该方法在 $n\leq6$ 传统函数上显著降低 T-count 与加权 score，在 ROS-style LUT proxy 上保持 309/0/0 的 score 优势，在 $n=18,20$ 函数级边界上验证了 \screen{} 与 Resource tail 的互补作用，在 $n=20$ 到 $40$ 项集级 scale 中展示了结构级策略的可扩展性，并通过 $n=21,22,23$ truth-table bridge 补强了高维正确性证据。Large frontier policy 将独立泛化集上相对 all-depth 的平均 score gap 压低到 0.10\%，并在 n=23 完整 bridge 上相对旧 policy 进一步降低 0.48\% score 和 0.45\% T-depth proxy。当前稿件已经形成独立于 SSHR 的投稿主线，但最终论文仍需官方 ROS/RevKit/mockturtle 对比和真正后端相关评估来提升审稿说服力。

\begin{thebibliography}{99}

\bibitem{meuli2019multiplicative}
G.~Meuli, M.~Soeken, E.~Campbell, M.~Roetteler, and G.~De Micheli.
\newblock The role of multiplicative complexity in compiling low T-count oracle circuits.
\newblock In \emph{Proceedings of ICCAD}, 2019.

\bibitem{meuli2022xag}
G.~Meuli, M.~Soeken, and G.~De Micheli.
\newblock Xor-And-Inverter Graphs for quantum compilation.
\newblock \emph{npj Quantum Information}, 8:7, 2022.

\bibitem{meuli2020ros}
G.~Meuli, M.~Soeken, M.~Roetteler, and G.~De Micheli.
\newblock ROS: Resource-constrained oracle synthesis for quantum computers.
\newblock \emph{Electronic Proceedings in Theoretical Computer Science}, 318:119--130, 2020.

\bibitem{wille2009bdd}
R.~Wille and R.~Drechsler.
\newblock BDD-based synthesis of reversible logic for large functions.
\newblock In \emph{Proceedings of DAC}, pp.~270--275, 2009.

\bibitem{yu2025backend}
M.~Yu, A.~Tempia Calvino, M.~Soeken, and G.~De Micheli.
\newblock Back-end-aware fault-tolerant quantum oracle synthesis.
\newblock In \emph{Proceedings of ASP-DAC}, pp.~930--937, 2025.

\bibitem{zheng2025sshr}
Q.~Zheng, Y.~Xu, J.~Zhang, Z.~Su, and S.~Zheng.
\newblock CNOT oriented synthesis for small-scale Boolean functions using spatial structures of parallelotopes.
\newblock In \emph{Proceedings of ICCAD}, 2025.

\bibitem{tsaras2024shortcircuit}
D.~Tsaras, A.~Grosnit, L.~Chen, Z.~Xie, H.~Bou-Ammar, and M.~Yuan.
\newblock ShortCircuit: AlphaZero-driven circuit design.
\newblock \emph{arXiv:2408.09858}, 2024.

\bibitem{rietsch2024unitary}
S.~Rietsch, A.~Y. Dubey, C.~Ufrecht, M.~Periyasamy, A.~Plinge, C.~Mutschler, and D.~D. Scherer.
\newblock Unitary synthesis of Clifford+T circuits with reinforcement learning.
\newblock In \emph{IEEE International Conference on Quantum Computing and Engineering}, pp.~824--835, 2024.

\bibitem{riu2025rlzx}
J.~Riu, J.~Nogu{\'e}, G.~Vilaplana, A.~Garcia-Saez, and M.~P. Estarellas.
\newblock Reinforcement learning based quantum circuit optimization via ZX-calculus.
\newblock \emph{Quantum}, 9:1758, 2025.

\bibitem{machiya2026monteq}
M.~Machiya, M.~Menickelly, P.~Hovland, and J.~Liu.
\newblock MonteQ: A Monte Carlo tree search based quantum circuit synthesis framework.
\newblock \emph{arXiv:2604.19029}, 2026.

\end{thebibliography}

\end{document}
